\begin{center}\begin{tikzpicture}[font=\small,>=stealth,line join=round]
\draw[gray!40,rounded corners=3pt] (0,0) rectangle (14.2,8.3);\node at (7.1,7.700000000000001) {例1：轴在1与3，0、2、4的高度都为7};
\begin{scope}[shift={(1.7774193548387096,1.88)},x=0.9677419354838709cm,y=0.45999999999999996cm]\draw[->,gray] (-0.7,0)--(5.5,0) node[right]{$x$};\draw[->,gray] (0,-0.5)--(0,9.5) node[above]{$y$};\draw[blue,thick,solid,domain=-0.3:5,samples=180] plot(\x,{7+cos(deg(pi*(\x-1)/2))});\fill[orange!85!black] (0,7) circle(2.2pt);\fill[orange!85!black] (2,7) circle(2.2pt);\fill[orange!85!black] (4,7) circle(2.2pt);\foreach \a in {1,3}{\draw[violet,dashed] (\a,5.5)--(\a,8.5);\node[above] at (\a,8.5){$x=\a$};}\end{scope}\node[anchor=west,align=left,text width=5.0cm] at (8.4,4.0) {$f=7+\cos\frac{\pi(x-1)}2$\\[9pt]$f(0)=f(2)=f(4)=7$\\[9pt]轴间距2，周期可取4\\[9pt]$2024=506\times4$\\[9pt]$f(2024)=7$};
\node[text width=12.6cm,align=center] at (7.1,.65) {图形与例题初值7相符；两次反射足以求值，无须确定函数解析式。};
\end{tikzpicture}\end{center}
